% =====================================================================
%  Indexing Without Prediction
%  Maximality and Empirical Boundary of the Lach Elemental Lattice
%  Matthew Lach
%
%  Compile:  pdflatex indexing-without-prediction.tex   (x2 for refs)
%  Requires: figs/fig01.png ... figs/fig10.png
% =====================================================================
\documentclass[11pt,a4paper]{article}

\usepackage[T1]{fontenc}
\usepackage[utf8]{inputenc}
\usepackage{lmodern}
\usepackage[margin=25mm]{geometry}
\usepackage{amsmath,amssymb,amsthm}
\usepackage{graphicx}
\usepackage{booktabs}
\usepackage{longtable}
\usepackage{array}
\usepackage{caption}
\usepackage{microtype}
\usepackage{enumitem}
\usepackage[dvipsnames]{xcolor}
\usepackage{fancyhdr}
\usepackage{titlesec}
\usepackage[hidelinks,pdfusetitle]{hyperref}

\definecolor{ink}{HTML}{1A1A1A}
\definecolor{acc}{HTML}{8C2D19}
\definecolor{acctwo}{HTML}{2F5D7C}
\definecolor{mut}{HTML}{6B6B6B}

\captionsetup{font=small,labelfont=bf,labelsep=period,justification=raggedright,
              singlelinecheck=false}

\pagestyle{fancy}
\fancyhf{}
\renewcommand{\headrulewidth}{0.4pt}
\fancyhead[L]{\footnotesize\color{mut}Indexing Without Prediction --- Lach Elemental Lattice}
\fancyfoot[C]{\footnotesize\color{mut}\thepage}
\fancypagestyle{plain}{\fancyhf{}\renewcommand{\headrulewidth}{0pt}%
  \fancyfoot[C]{\footnotesize\color{mut}\thepage}}

\titleformat{\section}{\normalfont\Large\bfseries\color{ink}}{\thesection.}{0.6em}{}
\titleformat{\subsection}{\normalfont\bfseries\color{ink}}{\thesubsection}{0.6em}{}
\titleformat{\subsubsection}{\normalfont\bfseries\color{acctwo}}{\thesubsubsection}{0.6em}{}

\newcommand{\parttitle}[1]{%
  \vspace{1.2\baselineskip}%
  {\noindent\LARGE\bfseries\color{acc} #1\par}%
  \vspace{0.4\baselineskip}\noindent\textcolor{acc}{\rule{\linewidth}{0.8pt}}%
  \vspace{0.6\baselineskip}\par}

\theoremstyle{definition}
\newtheorem{proposition}{Proposition}[section]
\newtheorem{lemma}[proposition]{Lemma}
\newtheorem{theorem}[proposition]{Theorem}
\newtheorem{corollary}[proposition]{Corollary}

\newcommand{\Lm}{\ensuremath{\Lambda}}
\newcommand{\Le}{\ensuremath{\Lambda_{8}}}
\newcommand{\Zeff}{\ensuremath{Z_{\mathrm{eff}}}}
\newcommand{\cmi}{\,\ensuremath{\mathrm{cm}^{-1}}}
\newcommand{\Jc}{\ensuremath{J_{\mathrm{c}}}}

\title{\vspace{-12mm}\textbf{\Huge Indexing Without Prediction}\\[3mm]
       {\large\color{acctwo} Maximality and Empirical Boundary of the
        Lach Elemental Lattice}}
\author{\textbf{Matthew Lach}\\[1mm]{\small\itshape Independent researcher}}
\date{}

\begin{document}

% =====================================================================
\begin{titlepage}
\thispagestyle{empty}
\vspace*{28mm}
{\noindent\Huge\bfseries Indexing Without Prediction\par}
\vspace{3mm}
{\noindent\large\color{acctwo} Maximality and Empirical Boundary of the Lach
 Elemental Lattice\par}
\vspace{4mm}
\noindent\textcolor{acc}{\rule{\linewidth}{1pt}}
\vspace{10mm}

{\noindent\large\textbf{Matthew Lach}\par}
{\noindent\small\color{mut} Independent researcher\par}
\vspace{10mm}

\noindent\textbf{An order-theoretic study of atomic configuration space, and a
certified-bound method for Rydberg energies.}
\vspace{3mm}

{\noindent\small\color{mut}
We construct a closed distributive lattice of order dimension exactly eight over
atomic electron configurations; prove that no quantity derived from its
coordinates can raise that dimension; show it extends by one dimension per
independent quantum number, reaching injectivity at eleven; and establish that
none of this yields predictive uplift. What does survive is a bracketing method
for Rydberg binding energies, derived entirely from the Rydberg formula, whose
deductive containment holds in 295 of 295 interior cells across six species,
three effective charges and six values of orbital angular momentum.\par}
\vspace{12mm}

\noindent\begin{tabular}{@{}p{58mm}p{95mm}@{}}
\toprule
Object & \Lm$_{8}$, extended to \Lm$_{11}$\\
Central negative result & order structure supplies indexing, not prediction\\
Central positive result & deductive containment, 295 / 295, no exceptions\\
Verified forward prediction & Kr\,\textsc{ii} $^2$P$^{\circ}_{1/2}$ splitting:
  $5591\pm500$ predicted, $5370.10$ measured\\
Claims withdrawn from earlier work & nine, listed in \S\ref{sec:withdrawn}\\
Errors found and corrected & seven, listed in Appendix~\ref{app:errors}\\
\bottomrule
\end{tabular}
\vspace{14mm}

{\noindent\small\color{mut}
\textbf{Status.} Prepared for independent review. Every numerical claim is
reproducible from the sources in Appendix~\ref{app:data}. Readers checking this
work are asked to attack \S\ref{sec:offset} first and \S\ref{sec:lemmas}
second; those are the two places where the argument is thinnest and the author
knows it.\par
\vspace{2mm}
\textbf{Data.} All levels are drawn from published critical compilations; no new
measurements are reported. Full provenance in Appendix~\ref{app:data}, full
bibliography in the References.\par}
\end{titlepage}

% =====================================================================
\section*{Abstract}
\addcontentsline{toc}{section}{Abstract}
\small
We construct a lattice \Le{} over atomic electron configurations, prove that no
quantity derived from its coordinates can raise its order dimension, and
establish empirically that its order structure supplies no predictive
information beyond the coordinates from which it is built.

\Le{} is a closed distributive lattice of order dimension exactly eight,
generated by an admissibility condition requiring every constraint to take the
form $x \le \varphi(y)$ with $\varphi$ non-decreasing in a single coordinate and
a constant floor. We prove that adjoining any function of the cell leaves the
dimension unchanged, because the graph of a monotone function is
order-isomorphic to its domain. The lattice is nonetheless extensible by
\emph{independent} physical quantities --- the parent ion's angular momentum,
the total angular momentum, the jK label, nuclear spin --- each contributing
exactly one dimension. \Lm$_{11}$ so constructed is injective on every test case
assembled here.

Against this we set a result of the opposite sign. A bracketing method for
Rydberg binding energies, derived entirely from $T=\Zeff^{2}R/\nu^{2}$ and using
none of the lattice's order structure, produces containment intervals holding in
295 of 295 interior cells across hydrogen, four alkaline earths, three
alkali-like ions and a noble gas --- spanning $\Zeff = 1$--3, $\ell = 0$--5,
single-, two- and multi-limit systems, bound and autoionising. Its uncertainty
estimate is derived rather than calibrated and returns $\pm1\sigma$ coverage of
68.3\% pooled over 293 cells, against a Gaussian expectation of 68.3\%.

Six independent attempts to extract predictive uplift from the lattice's order
all failed, each by a distinct mechanism. We show these reduce to one theorem:
on a product order, join, meet and comparability are definable from the
coordinates, so lattice quantities restate the coordinates but cannot exceed
them. The theorem is indifferent to dimension and to how faithfully the lattice
represents the physics. Over the course of this work the index was improved to
exactness; the predictive boundary did not move.

The one construction that escapes the theorem is deductive containment, because
a bracket is not a function of the coordinates but an entailment of the order
relation given monotonicity. That is the sense in which an ordinal structure
over configuration space is useful, and the only sense we have verified.
\normalsize

% =====================================================================
\parttitle{Part 0 --- Origin and Orientation}

\section{Introduction}

\subsection{Where this began}
This framework began as a three-dimensional recoding of the periodic table. Each
element was assigned a triple $(n,\ell,k)$: the principal quantum number of the
subshell being filled, its orbital angular momentum, and the occupancy of that
subshell. Under the componentwise partial order these triples form a lattice,
and the periodic table becomes a path through it.

The recoding is exact and it is not new physics --- it is a bijective
relabelling of information already present in the Madelung filling
order~\cite{madelung,klechkowski}. Its interest was structural: periodicity
appears as a projection of a three-dimensional object, and chemically analogous
elements are close in the lattice for a structural reason rather than by
convention~\cite{scerri,ostrovsky}.

Three additions produced the object studied here. Excited configurations require
distinguishing the subshell an electron leaves from the one it enters, giving a
second triple $(e,f,g)$. Ionisation requires a count $q$ of electrons removed.
Spin multiplicity requires a coordinate $S$. The result is \Le{}, whose cells
index electron configurations relative to a reference state.

The framework's earlier claims concerned prediction: that the lattice's order,
closure and dimension would constrain atomic properties in ways the coordinates
alone did not. This paper reports that those claims are false, establishes why
in a form we believe is a theorem rather than a limitation of effort, and
reports what survives.

\subsection{What this paper claims}
\textbf{Structural (Part I).} \Le{} is closed under join and meet and has order
dimension exactly eight. No derived quantity raises that dimension. It extends
by one dimension per independent physical degree of freedom. As an index it
achieves 86\% coverage and, at eleven dimensions, injectivity on every test
case. The physical subset is \emph{not} a sublattice, and no choice of
coordinates makes it one.

\textbf{Empirical (Part II).} A four-rule bracketing method, each rule derived
from $T=\Zeff^{2}R/\nu^{2}$, produces deductive containment verified on 295
interior cells with no exceptions, and an uncertainty estimate requiring no
calibration and returning 68.3\% coverage at $\pm1\sigma$. Its domain is bounded
by four exclusions, each with an identified mechanism.

\textbf{Thesis (Part III).} The lattice's order contributes nothing to the
method. Six tests establish this and a theorem about product orders explains
why. The theorem does not depend on dimension or on faithfulness, so improving
the index --- which Part I does --- leaves the predictive boundary where it was.

\subsection{What this paper does not claim}
We do not claim novel physics. Every quantitative result in Part II follows from
the Rydberg formula~\cite{rydberg} and the quantum defect~\cite{ritz,seaton};
the lattice's role is to make ``interior on some axis'' and ``same channel''
well defined. We do not claim the method extends beyond energies;
\S\ref{sec:domain} reports two attempts and both failed, for reasons we
characterise.

An earlier line of this work concerned applications to condensed-matter nuclear
phenomena. That material is removed entirely. It rested on extensions whose
structural status we now know to be defective (Appendix~\ref{app:multi}), and
revisiting it requires first solving the multi-target representation problem.

% =====================================================================
\parttitle{Part I --- Structure}

\section{Construction}

\subsection{The admissibility condition}
Let $\Lm \subseteq \mathbb{N}^{d}$ be defined by a family of constraints. The
family is \textbf{admissible} --- the condition previously designated \S9.4 ---
if every constraint has the form
\begin{equation}
x_{i} \;\le\; \varphi(x_{j}), \qquad \varphi \text{ non-decreasing},
\qquad\text{together with a constant floor } x_{i}\ge c .
\label{eq:admis}
\end{equation}

\begin{proposition}
An admissible constraint family defines a set closed under componentwise join
and meet.
\end{proposition}
\begin{proof}
Let $a,b\in\Lm$. For the join, $\max(a_i,b_i)$ equals one of them, say $a_i$;
then $a_i \le \varphi(a_j) \le \varphi(\max(a_j,b_j))$ since $\varphi$ is
non-decreasing. For the meet, choose the index $m\in\{a,b\}$ minimising $x_j$;
then $\min(a_i,b_i)\le m_i \le \varphi(m_j)=\varphi(\min(a_j,b_j))$. A constant
floor is preserved trivially.
\end{proof}

Two consequences shape everything following. A bound on a \textbf{sum} of two
coordinates is not admissible and generally breaks closure --- this defeats the
multi-target extension of Appendix~\ref{app:multi}. A \textbf{varying floor} is
likewise inadmissible, which places the physical selection rules outside the
language (\S\ref{sec:notsub}).

\subsection{The eight axes}
\begin{center}\small
\begin{tabular}{@{}llll@{}}
\toprule
axis & meaning & bound & form\\
\midrule
$n$    & source shell          & $1\le n\le N$              & free root\\
$\ell$ & source subshell       & $\ell \le n-1$             & non-decr.\ in $n$\\
$k$    & source occupancy      & $1\le k\le 2(2\ell+1)$     & non-decr.\ in $\ell$\\
$q$    & electrons removed     & $0\le q\le k$              & non-decr.\ in $k$\\
$e$    & target shell          & $1\le e\le E$              & free root\\
$f$    & target subshell       & $f\le e-1$                 & non-decr.\ in $e$\\
$g$    & electrons transferred & $g\le\min(q,2(2f+1))$      & min of two admissible\\
$2S$   & spin multiplicity     & $0\le 2S\le k$             & non-decr.\ in $k$\\
\bottomrule
\end{tabular}
\end{center}
\vspace{2mm}

Derived quantities include net charge $c=q-g$, electron count $N$, and target
occupancy $G$. \textbf{These are derived precisely because they cannot be
coordinates}: each is a difference or sum of two axes, hence inadmissible.
\S\ref{sec:rejected} makes this exact.

\subsection{Closure and dimension}
Closure follows from Proposition~2.1 and was verified additionally on
$8\times10^{6}$ sampled pairs at caps covering the observed data, with zero
failures.

\textbf{Order dimension is exactly eight.} The upper bound is immediate. The
lower bound follows from an explicit 8-box at
$(n,\ell,k,q,e,f,g,2S)=(3,1,2,1,2,0,0,0)$, all $2^{8}=256$ corners of which are
admissible. The box is \textbf{cap-independent}, remaining valid at cap settings
$(4,2,4,2)$, $(6,3,8,3)$ and $(8,4,12,4)$. Order dimension is used throughout in
the sense of Dushnik and Miller~\cite{dushnik}; see Trotter~\cite{trotter} for
the standard results invoked.

\begin{figure}[htbp]\centering
\includegraphics[width=0.86\linewidth]{figs/fig01.png}
\caption{The admissibility constraint graph. Circles are axes; an arrow
$y\to x$ indicates $x$ is bounded by $\varphi(y)$. Two roots ($n$ and $e$) make
the graph a forest with a single cross-link at $g$.}
\end{figure}

\begin{figure}[htbp]\centering
\includegraphics[width=\linewidth]{figs/fig02.png}
\caption{The 8-box projected onto four coordinate pairs. Every one of the 256
corners is admissible; labels give the number of corners mapping to each
projected point.}
\end{figure}

\subsection{A note on cell counts}
Earlier versions reported $|\Le| = 61\,453$ as a structural invariant. It is
not. That is the cell count at one choice of caps, and it changes with them.
Every quantity derived from it --- the fraction of cells representing inward
transitions, the fraction ``realisable'' --- is a truncation artefact and is
withdrawn. The invariants of \Le{} are its closure, its order dimension and its
behaviour under extension; its cardinality is not among them.

\section{Algebra and geometry, axis by axis}
Each axis has an \emph{algebra} --- the bound defining it and its admissible
form --- and a \emph{geometry} --- the shape of the region it occupies. This
section also reports where each is exercised, which varies enormously.

\subsubsection*{$n$, the source shell}
Algebra: a free root; every electronic bound descends from it. Geometry: a
chain. Exercised: distinguishes species but is \textbf{constant within any one
atom}. Removing it from a single species' index costs zero collisions.

\subsubsection*{$\ell$, the source subshell}
Algebra: $\ell\le n-1$, non-decreasing in $n$. Geometry: a right triangle in the
$(n,\ell)$ plane. Exercised: distinguishes species and blocks; constant within
an atom.

\subsubsection*{$k$, the source occupancy}
Algebra: $1\le k\le 2(2\ell+1)$. Geometry: a staircase widening with $\ell$ ---
2, 6, 10, 14. Exercised: constant within an atom. $k$ bounds $q$, $2S$ and the
angular-momentum axes, so it is structurally central while empirically inert
within a species.

\subsubsection*{$q$, electrons removed}
Algebra: $0\le q\le k$, a triangle nested inside $k$. Exercised: separates
neutral from ionised species; constant ($q=1$) across all singly-excited states
of one atom.

\subsubsection*{$e$, the target shell}
Algebra: a second free root. \textbf{This is the workhorse.} Removing $e$ from
the Ca\,\textsc{i} index costs 13 collisions in 27 levels; from
Kr\,\textsc{i}, 22 in 46. It is the axis the method of Part II interpolates
along, and the only axis load-bearing both across and within species.

\subsubsection*{$f$, the target subshell}
Algebra: $f\le e-1$, a second right triangle independent of the $(n,\ell)$ one.
Geometry: \Le{} is therefore a product of two triangles, not a simplex --- the
clearest sense in which it is not a relabelled periodic table. Exercised: 7
collisions for Ca\,\textsc{i}, 22 for Kr\,\textsc{i}.

\subsubsection*{$g$, electrons transferred}
Algebra: $g\le\min(q,2(2f+1))$ --- the intersection of two triangles, a wedge,
and the only axis bounded by two others. Exercised: \textbf{zero collisions on
removal.} For singly-excited states $g=q=1$ universally, so the axis is
exercised only by doubly-excited configurations, of which the record holds few.

\subsubsection*{$2S$, spin}
Algebra: $0\le 2S\le k$, a chain over $k$. Exercised: 5 collisions for
Ca\,\textsc{i}. Load-bearing where LS coupling holds; \textbf{undefined} for
jK-coupled systems such as Kr\,\textsc{i}, where the parent's $J$ must
substitute.

\subsection{The collective whole}
\Le{} is a product of two triangular regions --- $(n,\ell)$ and $(e,f)$ ---
with occupancy chains hung from each and a wedge ($g$) linking them. It is not a
simplex, not a hypercube and not a chain product; two roots make the constraint
graph a forest with a single cross-link.

\begin{figure}[htbp]\centering
\includegraphics[width=0.92\linewidth]{figs/fig03.png}
\caption{The collective whole. Two triangular blocks, two chains descending from
$k$, and one cross-link where $g$ is bounded by both $q$ and $f$.}
\end{figure}

\subsection{Axis ablation}
Removing each axis in turn from the index and counting collisions measures
directly what each contributes.

\begin{center}\small
\begin{tabular}{@{}lcc@{}}
\toprule
axis & Ca\,\textsc{i} (27 levels) & Kr\,\textsc{i} (46 levels)\\
\midrule
$n,\ell,k,q,g$ & 0 & 0\\
$e$            & 13 & 22\\
$f$            & 7  & 22\\
$2S$           & 5  & undefined (jK)\\
$2\Jc$ (core)  & 0 (fixed) & 20\\
$2J$ (total)   & 12 & 35\\
\bottomrule
\end{tabular}
\end{center}
\vspace{2mm}

Six of ten coordinates are constant within a species and discriminate only
\emph{across} species. Only $e$ and $f$ appear in both roles --- which is
exactly why the method of Part II interpolates in $e$ at fixed $f$ and uses no
other axis.

\begin{figure}[htbp]\centering
\includegraphics[width=0.92\linewidth]{figs/fig04.png}
\caption{Axis ablation for an LS-coupled species (Ca\,\textsc{i}) and a
jK-coupled one (Kr\,\textsc{i}). Six of ten coordinates are inert within a
species.}
\end{figure}

\section{What cannot be added, and what can}
\label{sec:lemmas}

\subsection{Two lemmas and a corollary}
The result is best stated as two facts, because they carry unequal weight and
the combined form disguises which is which.

\begin{lemma}[the substantive half]
Let $L$ be a finite lattice and $h:L\to\mathbb{Z}$ a function whose graph
$G=\{(x,h(x))\}$ is closed under componentwise join and meet. Then $h$ is
monotone.
\end{lemma}
\begin{proof}
Closure means that for any $x,y\in L$ the pair $(x\vee y,\max(h(x),h(y)))$ lies
in $G$. Since $G$ is the graph of a function its second component must equal $h$
of its first, so $h(x\vee y)=\max(h(x),h(y))$. If $x\le y$ then $x\vee y=y$,
giving $h(y)=\max(h(x),h(y))\ge h(x)$.
\end{proof}

\begin{lemma}[bookkeeping]
If $h:L\to\mathbb{Z}$ is monotone then $x\mapsto(x,h(x))$ is an order
isomorphism of $L$ onto $G$, and hence $\dim(G)=\dim(L)$.
\end{lemma}
\begin{proof}
The map is a bijection onto $G$. Monotonicity makes it order-preserving;
projection onto the first coordinate makes it order-reflecting.
\end{proof}

\begin{theorem}
Adjoining to $L$ any function of its cells whose graph is closed leaves the
order dimension unchanged.
\end{theorem}
\begin{proof}
Immediate from Lemmas 4.1 and 4.2.
\end{proof}

The division of labour is worth stating plainly. Lemma~4.2 requires only
monotonicity and would hold for any monotone function whatever; it is a remark
about graphs, not a result about lattices. All the content sits in Lemma~4.1 ---
that closure is strong enough to force monotonicity in the first place.
Presenting the two as a single theorem, as an earlier version of this paper did,
makes the conclusion look deeper than it is. Neither lemma needs an
admissibility hypothesis, and neither assumes $L$ is a product of chains.

\begin{corollary}
No quantity computed from the eight electronic coordinates can raise \Le{}'s
order dimension.
\end{corollary}

\paragraph{Independent verification.} Four checks were run against this
argument.
\begin{enumerate}[leftmargin=*,itemsep=1pt]
\item The $d$-box test used throughout the paper is validated: exact order
dimension computed from minimum realizers returns 2 for $2^{2}$ and 3 for
$2^{3}$, and since a $d$-box is an induced copy of $2^{d}$ while order dimension
is monotone under induced subposets~\cite{trotter}, containing one forces
$\dim\ge d$.
\item Exhaustively over all $19\,683$ functions from a $3\times3$ grid into a
three-value codomain, \textbf{no} non-monotone function has a closed graph ---
confirming Lemma~4.1.
\item Across 175 monotone functions there are zero mismatches between the order
on $L$ and the order on its graph --- confirming Lemma~4.2.
\item On a lattice carrying the $\ell\le n-1$ constraint that breaks
interval-closure, exact realizer-based dimension is 2 for the base and 2 for all
ten distinct closure-preserving extensions, confirmed against true dimension
rather than the box lower bound. Exact dimension is computed only for small
lattices, the problem being NP-hard in general~\cite{yannakakis}.
\end{enumerate}

\paragraph{A stronger claim that is false.} An earlier version of this work
asserted that any closure-preserving $h$ must be a monotone function of a
\emph{single} coordinate. That holds for products of chains but not for \Lm{},
which is a sublattice of one: the constraint $\ell\le n-1$ means \Lm{} is not
interval-closed, since $(1,2)$ lies componentwise between $(1,0)$ and $(3,2)$
yet is not a cell. On the test lattice, 148 of 335 closure-preserving functions
depend on more than one coordinate --- for instance the one generated by the
chain $(1,0,2)\le(3,0,2)$, which varies with both $n$ and $\ell$. Those
functions are nonetheless redundant, not because they depend on one axis but
because they add no dimension.

\subsection{Rejected candidates}
\label{sec:rejected}
Each of the following was proposed as a ninth axis in earlier work; each fails,
and the failure is now attributable to Theorem~4.3 rather than to case-by-case
testing.

\begin{center}\small
\begin{tabular}{@{}llrrl@{}}
\toprule
candidate & form & join fail & meet fail & verdict\\
\midrule
$h=n$                  & projection      & 0       & 0       & closed, no 9-box\\
$h=2n$                 & monotone in $n$ & 0       & 0       & closed, redundant\\
$h=n+\ell$             & two coordinates & 42\,492 & 42\,492 & breaks closure\\
$h=q-g$ (net charge)   & difference      & 284\,858& 284\,858& breaks closure\\
$h=k-q$                & difference      & 406\,400& 406\,400& breaks closure\\
$h=e-f$                & difference      & 266\,440& 266\,440& breaks closure\\
$h=\lfloor C/e^{2}\rfloor$ & anti-monotone & 503\,404 & 503\,404 & breaks closure\\
\bottomrule
\end{tabular}
\end{center}
\vspace{2mm}

The energy proxy deserves separate comment, being the original motivation for a
ninth ``output'' axis. It fails for a sharper reason than inadmissibility:
energy \emph{decreases} with $e$ while the order increases, so an energy-valued
coordinate is anti-monotone and pairs of cells acquire no upper bound at all.
\textbf{A metric quantity increasing with the coordinates might well be
admissible.} The obstruction is the sign, not the metricity.

\section{Extension by independent quantities}

\subsection{Three classes}
Theorem~4.3 constrains only functions of the existing coordinates. Physically
independent quantities escape it, and fall into three classes.
\begin{enumerate}[label=(\alph*),leftmargin=*,itemsep=1pt]
\item \textbf{Functions of the cell.} Graph order-isomorphic to \Lm{}; dimension
unchanged. Always.
\item \textbf{Independent, bounded by an existing coordinate.} Admissible if and
only if the bound has form~\eqref{eq:admis}. Contributes exactly one dimension.
\item \textbf{Independent, bounded by a constant.} The set is a direct product;
closure automatic; one dimension.
\end{enumerate}

\subsection{$\Lambda_{9}$, $\Lambda_{10}$, $\Lambda_{11}$}
\textbf{Parent-ion angular momentum (class b).} The same electronic cell can
leave the residual ion in different $J$ states, so \Jc{} is not a function of
the eight. With $2\Jc\le k$ --- non-decreasing in a single coordinate --- the
enlarged set is \textbf{closed exactly}: 6989 cells, all $24\,419\,566$ pairs
tested, zero failures, with a 9-box at $(3,1,5,4,3,1,3,2,1)$. Order dimension 9.

\textbf{Total angular momentum (class b).} With $2J\le2k$, closure holds on
$2\,041\,470$ cells with zero failures and a 10-box at
$(3,1,2,1,2,0,0,0,0,0)$. Order dimension 10.

\textbf{jK label (class b)}, required for jK-coupled
systems~\cite{cowan,condonshortley}, and \textbf{nuclear spin (class c)},
bounded by a constant, follow the same pattern.

\subsection{What extension buys and costs}
\begin{center}\small
\begin{tabular}{@{}lccl@{}}
\toprule
lattice & Kr\,\textsc{i} cells & collisions & Ca\,\textsc{i} fine structure\\
\midrule
$\Lambda_{8}$  & 7  & 39 & 1 collision\\
$\Lambda_{9}$  & 11 & 35 & 1 collision\\
$\Lambda_{10}$ & 34 & 12 & resolved\\
$\Lambda_{11}$ & 46 & 0  & resolved\\
\bottomrule
\end{tabular}
\end{center}
\vspace{2mm}

The cost is looseness. Each class-(b) bound admits more than the physics
permits: the true range for total $J$ is the triangle rule
$|\Jc-j|\le J\le \Jc+j$, whose \textbf{lower bound varies with the cell}, and
\eqref{eq:admis} permits only constant floors. On a bounded enumeration the
physically realisable fraction is 23.3\% at eight dimensions, 9.0\% at nine,
1.7\% at ten.

\begin{figure}[htbp]\centering
\includegraphics[width=0.86\linewidth]{figs/fig05.png}
\caption{Extension clears collisions (left axis) while admitting unphysical
cells (right axis). The two curves crossing is the trade.}
\end{figure}

\section{The lattice as an index}

\subsection{Coverage}
\begin{center}\small
\begin{tabular}{@{}lcc@{}}
\toprule
species & mapped & coverage\\
\midrule
Kr\,\textsc{i}   & 9 / 9   & 100\%\\
Sc\,\textsc{iii} & 8 / 8   & 100\%\\
Ba\,\textsc{ii}  & 10 / 10 & 100\%\\
Sr\,\textsc{i}   & 9 / 10  & 90\%\\
Ca\,\textsc{i}   & 12 / 15 & 80\%\\
Ti\,\textsc{i}   & 6 / 11  & 55\%\\
\midrule
total            & 54 / 63 & 86\%\\
\bottomrule
\end{tabular}
\end{center}
\vspace{2mm}

Every failure is a \textbf{multi-subshell configuration} --- two occupied
non-source subshells, which a single $(e,f,g)$ triple cannot name. Coverage is a
representational limit and \textbf{no additional coordinate repairs it}; the
natural repair, a second target triple, breaks the lattice
(Appendix~\ref{app:multi}).

\subsection{Injectivity}
\begin{center}\small
\begin{tabular}{@{}llll@{}}
\toprule
collision type & example & resolved by & at\\
\midrule
core $J$       & Kr\,\textsc{i} $(^2$P$^{\circ}_{3/2})5s$ vs
                 $(^2$P$^{\circ}_{1/2})5s$ & $2\Jc$ & $\Lambda_{9}$\\
fine structure & Ca\,\textsc{i} $4s4d\ ^3$D, $J=1,2,3$ & $2J$ & $\Lambda_{10}$\\
jK label       & Kr\,\textsc{i} $5p\ ^2[1/2]$ vs $^2[3/2]$, same $J$
               & $K$ & $\Lambda_{11}$\\
\bottomrule
\end{tabular}
\end{center}
\vspace{2mm}

The first is consequential beyond bookkeeping. The two Kr\,\textsc{i} channels
converge on limits \textbf{5370\cmi{} apart}, and the method of Part II requires
knowing which. \textbf{\Le{} cannot supply that assignment}; in the work
reported here it was read from the ASD configuration string~\cite{nistasd}.
$\Lambda_{9}$ supplies it.

\subsection{The physical set is not a sublattice}
\label{sec:notsub}
The physically realisable cells satisfy the spin selection
rules~\cite{condonshortley}: $2S\equiv k \pmod 2$ and
$2S\le\min(k,2(2\ell+1)-k)$. \textbf{This set is not closed.} In the $(k,2S)$
plane, $(1,1)$ and $(2,0)$ are both physical; their join $(2,1)$ is not, since
$2S$ must be even when $k$ is even. Failures: 1 at $\ell=0$, 15 at $\ell=1$, 70
at $\ell=2$.

The obstruction is structural. A \textbf{congruence is not an inequality}, and
\eqref{eq:admis} admits only inequalities with constant floors. The selection
rules of atomic physics --- spin parity, hole symmetry, the triangle rule ---
are none of them of that form.

\subsection{Offset coordinates: faithful but not closed}
\label{sec:offset}
\emph{This section reports a negative result obtained late in the work and
corrects an intermediate claim. It is the place a reviewer should attack first.}

One may attempt to bring the selection rules inside the language by
reparameterisation. Replace $2S$ by the paired-unit count $m=(k-2S)/2$, so that
$2S=k-2m$ carries the parity of $k$ automatically; replace $2J$ by the triangle
offset $t=(2J-|2\Jc-2j|)/2$, so the triangle rule holds by construction. Each
transformation is bijective, so the dimension is unchanged and Lemma~4.2 is
untouched.

Tested \textbf{in isolation}, each works: the $(k,m)$ plane is closed at every
$\ell$, and the $(\Jc,j,t)$ block is closed with 140 cells and zero failures.

\textbf{Assembled into one lattice, they do not.} Exhaustive test on 1431 cells
and all $1\,023\,165$ pairs: \textbf{$222\,075$ join failures and $11\,696$ meet
failures}. The cause is visible in the bounds --- $m$ has a floor varying with
$k$, and $2\Jc\le k-q$ is decreasing in $q$. Isolated closure does not imply
joint closure.

What the offset coordinates do achieve is worth recording: all 1431 cells decode
to genuine physical states, zero violations, and the encoding is injective ---
1903 distinct states, no duplicates. \textbf{Offset coordinates are a faithful
and injective parameterisation of the physical set that is not a lattice.}

The section's result is therefore: \textbf{one may have closure or faithfulness,
not both.} The lattice properly contains the physics; the physics is not a
sublattice; and no relabelling repairs this, because the obstruction is the
varying floor and not the choice of origin.

% =====================================================================
\parttitle{Part II --- A Bracketing Method for Rydberg Energies}

\section{Derivation}
Every rule below follows from one expression,
\begin{equation}
T \;=\; \frac{\Zeff^{2}R}{\nu^{2}}, \qquad \nu \;=\; n-\delta(n,\ell,\text{core}),
\label{eq:rydberg}
\end{equation}
where $T$ is the binding energy relative to the series limit and $\delta$ the
quantum defect~\cite{rydberg,ritz,seaton}. The method uses the lattice only to
make two notions well defined: which cells belong to the same \textbf{channel},
and which cells are \textbf{interior} on an axis.

\section{The four rules}
\begin{quote}\small
\textbf{Rule 1 --- Interiority.} Predict only cells with measured neighbours on
both sides along some axis, and only cells belonging to the same Rydberg
sequence as those neighbours.

\textbf{Rule 2 --- Variable.} Read the interpolation variable off
\eqref{eq:rydberg} for the axis in question: $E$ along $n$ within a series;
interval ratios along $J$; geometric in $T$ along the core shell.

\textbf{Rule 3 --- Separation.} Report the deductive bracket separately from the
inferential point estimate. Given monotonicity the true value cannot lie outside
$[\max\text{ below},\ \min\text{ above}]$.

\textbf{Rule 4 --- Uncertainty.} Fit $\delta$ locally by a Ritz
expansion~\cite{ritz} $\delta=\delta_{0}+\delta_{2}/n^{2}$ and take
$\sigma = 2R\,\Zeff^{2}\,\mathrm{SE}_{\mathrm{pred}}/\nu^{3}$.
\end{quote}

\textbf{Rule 1's second clause is not redundant.} Sc\,\textsc{iii} $3d$ is
interior in $n$ within $\ell=2$, yet including it inflates the median error by
$3.2\times$ because it is the collapsed ground state rather than a sequence
member. The operational test is that its defect departs from its neighbours' by
0.14 where they depart from each other by 0.008 --- a factor of 17.

Each rule is justified by the cost of violating it:

\begin{center}\small
\begin{tabular}{@{}ll@{}}
\toprule
rule violated & cost\\
\midrule
extrapolate rather than interpolate & rms $248\to928$\cmi{} ($3.7\times$), $n=91$\\
wrong interpolation variable        & rms $248\to589$\cmi{} ($2.4\times$)\\
pooled $\sigma$ instead of per-cell & $\pm1\sigma$ coverage 96.2\% instead of 68\%\\
include collapsed orbital           & median error $12.8\to41.3$\cmi{} ($3.2\times$)\\
\bottomrule
\end{tabular}
\end{center}

\section{Verification}
\textbf{295 interior cells, bracket exact in every one.} Two counts are reported
throughout and they differ by two: a bracket requires one measured neighbour on
each side, so 295 cells have one; the local Ritz fit has two free parameters and
so requires three neighbours, giving 293 cells with a derived $\sigma$. Pooled
over all 293, $\pm1\sigma$ coverage is \textbf{68.3\%} and $\pm2\sigma$ is
91.8\%, against Gaussian expectations of 68.3\% and 95.4\%.

\begin{center}\small
\begin{tabular}{@{}lcccr@{}}
\toprule
regime & cells & bracket & $\pm1\sigma$ & median error\\
\midrule
H\,\textsc{i} ($\Zeff=1$)                  & 5   & 5 / 5     & ---    & 0.96\cmi\\
Mg/Ca/Sr/Ba\,\textsc{i} (neutral)          & 189 & 189 / 189 & 67.7\% & 0.69\\
Ca/Sr/Ba\,\textsc{ii} (alkali-like)        & 37  & 37 / 37   & 70.3\% & 11.5\\
Ca\,\textsc{i} $3d$ channel (autoionising) & 22  & 22 / 22   & 68.2\% & 9.37\\
Kr\,\textsc{i} (multi-limit p-block)       & 31  & 31 / 31   & 76.7\% & 0.96\\
Sc\,\textsc{iii} ($\Zeff=3$)               & 11  & 11 / 11   & 54.5\% & 12.8\\
\bottomrule
\end{tabular}
\end{center}
\vspace{2mm}

\textbf{On independence.} These 295 cells are drawn from approximately
\textbf{22 channels}. Cells within a channel share neighbours, so successive
brackets are correlated and 295 should not be read as 295 independent tests.

\textbf{On the bracket's robustness.} Perturbing the series limit by
$\pm5$\cmi{} leaves coverage unchanged at 15/15, because the bracket uses only
measured neighbours. The point estimate and $\sigma$ require the limit;
\textbf{the bracket does not}. Similarly, ``interior'' requires that neighbours
exist, not that they be adjacent: with $n=8$--16 removed from a Ba\,\textsc{i}
series, cell 17 still brackets correctly, at width 2012\cmi{} instead of 400.

\textbf{On the neighbour count.} The Ritz fit uses the $k$ nearest members.
Sensitivity: rms 14.0 / 17.1 / 15.8 / 15.5 / 15.5\cmi{} for $k=3,4,6,8,10$. No
strong preference; $k=6$ was chosen arbitrarily and the sensitivity is reported
rather than the choice being presented as derived.

\begin{figure}[htbp]\centering
\includegraphics[width=0.88\linewidth]{figs/fig06.png}
\caption{Quantum-defect sequences across four species, two charge states and
$\ell=2,3$. The smooth Ritz behaviour is what Rule 4 fits and what the bracket
exploits.}
\end{figure}

\begin{figure}[htbp]\centering
\includegraphics[width=0.70\linewidth]{figs/fig07.png}
\caption{Derived $\sigma$ against realised error for all 293 cells carrying a
$\sigma$. The calibration is not fitted: $\sigma$ comes from the prediction
standard error of the local Ritz fit, and 68.3\% of cells fall within $1\sigma$
against a Gaussian expectation of 68.3\%.}
\end{figure}

\begin{figure}[htbp]\centering
\includegraphics[width=0.88\linewidth]{figs/fig08.png}
\caption{The bracket construction. One cell is withheld; the interval formed by
its measured neighbours must contain it, given monotonicity in $n$.}
\end{figure}

\section{Domain and exclusions}
\label{sec:domain}
\textbf{Verified domain.} Binding energies of Rydberg series: one electron
outside a core in a definite state, channel by channel, with the channel's limit
supplied by the configuration label. Interior cells only, excluding collapsed
orbitals.

\paragraph{(i) Sign-changing quantities.} Any property built from a matrix
element passing through zero cannot be monotone. Sr\,\textsc{i}
$5s^{2}\,^1$S$_0\to5snp\,^1$P$_1$ reduced transition probabilities collapse by a
\textbf{factor of 60} at $n=6$ and recover --- a Cooper-type
node~\cite{cooper}. Bracket coverage across the node: 79\%. Away from it: 100\%
on 11 cells.

\paragraph{(ii) Asymptotically constant quantities.} Fine-structure splittings
reduced by $\nu^{3}$ are constant to ${\sim}0.5\%$ across a series, and the
adjacent steps fall below measurement noise (1.6\%). Bracket coverage:
\textbf{35.3\%}; zero of 17 adjacent steps are resolvable. A constant is not
monotone.

\paragraph{(iii) Systems without interior cells.} Ti\,\textsc{i}: 23 channels,
26 members, \textbf{zero interior cells}. No channel reaches three members,
because the ionisation limit near $55\,000$\cmi{} lies below the onset of severe
level density, and above ${\sim}42\,000$\cmi{} the levels are 40--50\%
configuration-mixed with no assignable parent. The method declines rather than
errs.

\paragraph{(iv) Collapsed orbitals}, excluded by the defect-departure test of
Rule 1.

\begin{figure}[htbp]\centering
\includegraphics[width=\linewidth]{figs/fig09.png}
\caption{The three exclusion mechanisms, each measured: a sign change in the
matrix element; a quantity too flat to bracket; a system whose series are too
short to contain an interior cell.}
\end{figure}

\begin{figure}[htbp]\centering
\includegraphics[width=0.86\linewidth]{figs/fig10.png}
\caption{The verified domain. Cell counts by effective charge and orbital
angular momentum; the excluded region is marked with its mechanism.}
\end{figure}

\section{Predictions}

\subsection{Verified}
\textbf{Kr\,\textsc{ii} $^2$P$^{\circ}_{1/2}$ fine-structure splitting.} Derived
from three $\ell$ channels of the Kr\,\textsc{i} $(^2$P$^{\circ}_{1/2})$ series
and committed \emph{before} the measurement entered the analysis:
\textbf{$5591\pm500$\cmi}. Measured: \textbf{$5370.10$\cmi}~\cite{nistasd}.
Error $+221$\cmi, 4.1\%, within the stated interval.

\emph{Full disclosure on the interval.} The three channel estimates were 6061,
5679 and 5031, giving $\mathrm{sd}=521$. The quoted $\pm500$ is therefore
approximately $1\sigma$, not the conventional $2\sigma$ (which would be
$\pm1030$). The measured value lies $0.4\sigma$ from the mean. The prediction is
correct; the interval was tighter than convention and this should have been
stated at the time.

\subsection{Open}
\textbf{Ba\,\textsc{i} $5d^{2}\,^1$G$_4$.} Deductive bracket from the term
ordering $^3$F $<\ ^1$D $<\ ^3$P $<\ ^1$G $<\ ^1$S~\cite{racah}:
\textbf{$[23\,694,\ 26\,757]$\cmi}. No point estimate is offered; see
\S\ref{sec:withdrawn}.

\begin{center}\small
\begin{tabular}{@{}lcc@{}}
\toprule
Ba\,\textsc{i} level & bracket (\cmi) & Ritz estimate\\
\midrule
$6s8f$  & $[39\,678,\ 40\,614]$ & $40\,238$\\
$6s13f$ & $[41\,251,\ 41\,648]$ & $41\,367$\\
$6s14f$ & $[41\,251,\ 41\,648]$ & $41\,461$\\
$6s15f$ & $[41\,251,\ 41\,648]$ & $41\,536$\\
$6s16f$ & $[41\,251,\ 41\,648]$ & $41\,597$\\
\bottomrule
\end{tabular}
\end{center}
\vspace{2mm}

Leave-one-out validation on the same series: 15/15 bracket coverage, rms
20.5\cmi, with sub-wavenumber accuracy for $n\ge17$.

\textbf{Ti\,\textsc{ii} $a\,^4$F intervals}, derived from the Ti\,\textsc{i}
$4f$ and $5g$ manifolds resolved by core $J$, which agree to 0.7\cmi:
\textbf{94.3, 131.9, 167.3\cmi} for $J=3/2\to5/2\to7/2\to9/2$. Corresponding
limits: $55\,044.3$ / $55\,138.5$ / $55\,270.0$ / $55\,437.4$\cmi. Internal
check: taking $\delta(5g)=0$ gives $\delta(4f)=0.0508$, 0.0508, 0.0506, 0.0506
across four independent channels.

% =====================================================================
\parttitle{Part III --- Indexing Without Prediction}

\section{The theorem}
\begin{theorem}
On a product order, join, meet and comparability are definable from the
coordinates: $(a\vee b)_i=\max(a_i,b_i)$, $(a\wedge b)_i=\min(a_i,b_i)$, and
$a\le b \iff a_i\le b_i$ for all $i$.
\end{theorem}

Every lattice-theoretic quantity is therefore a function of the coordinates, and
can restate them but cannot exceed them. Two properties of this theorem drive
Part III. It is \textbf{indifferent to dimension}: $\Lambda_{11}$ is as
constrained as \Le{}. And it is \textbf{indifferent to faithfulness}: perfecting
the index moves nothing.

\section{Six mechanisms}

\subsubsection*{Rank aggregation}
Using $\mathrm{rank}(a\wedge b)$ and $\mathrm{rank}(a\vee b)$ as predictors of
configuration mixing: leave-one-out rmse 1.643 against 0.812 for a linear model
in $\Delta q,\Delta f$ --- worse than predicting the mean (data sd 0.940).
\emph{Mechanism:} rank is a \textbf{sum over axes}, and summing discards which
axis moved. Mixing depends on which subshell is involved; rank cannot see it.
This holds at any faithfulness.

\subsubsection*{Comparability}
Comparable cells order energies correctly in 65 of 65 Ca\,\textsc{i} pairs, and
select a 52.1\% subset with 100\% accuracy using no training data.
\emph{Mechanism:} energy is separately monotone in every coordinate --- 49
single-coordinate pairs, zero violations --- so $a\le b \Rightarrow E(a)\le E(b)$
is a \textbf{theorem given monotonicity}, not an empirical finding. Monotonicity
is a coordinate property that a positive-weight linear model encodes directly,
which is why the control reached 99.9\%.

\subsubsection*{Zero-shot subset selection}
The lattice selects a reliable subset with no labels; a fitted linear model
needs 40 labels to approach it. \emph{Mechanism:} a \textbf{naive
positive-weight sum}, all weights $+1$ and no fitting, matches it exactly ---
100\% at identical coverage, zero labels.

\subsubsection*{Closure as existence}
Observed configurations are closed under meet more than chance allows ---
Mg\,\textsc{i} $p=0.006$, Ca\,\textsc{i} 0.031, Sr\,\textsc{i} and
Ba\,\textsc{i} $<0.0001$ against a marginal-matched null. \emph{Mechanism:}
spectroscopists measure \textbf{downward-closed regions}. Every Rydberg series
is followed from its lowest member upward, so the observed set is a lower set by
construction.

\subsubsection*{Join-based configuration mixing}
For divalent atoms roughly 62\% of pairs have \textbf{no upper bound at all}, at
every valence count tested ($k=2,3,4$). \emph{Mechanism:} joining two
doubly-excited configurations requires transferring the union of their target
electrons, but $q\le k$. The join is not ambiguous; it does not exist.

\subsubsection*{A derived ordinal coordinate}
Collapsing four verified monotonicities into
$\kappa=(q,n_{\mathrm{core}},-n,-\ell)$ gives 2163/2163 monotone, 31.2\% of
pairs comparable, and containment across species and charge simultaneously.
\emph{Mechanism:} every component is a coordinate; the construction restates
Theorem~12.1 rather than escaping it. Its one apparent violation traced to a
transcription error in the input data (Appendix~\ref{app:errors}).

\section{What escapes}
Deductive containment escapes Theorem~12.1, and it is the only thing we have
found that does.

A bracket is not a function of the coordinates. It is an \textbf{entailment}:
given that the property is monotone in the order, the true value at an
unmeasured cell \emph{cannot} lie outside the interval formed by its measured
neighbours. A regression returns a point and a confidence interval that fails
some fraction of the time; the order returns containment that holds by
construction. That is why the bracket never failed in 295 cells while every
point-estimate claim in this work required qualification.

This also explains an asymmetry in compositional prediction. Propagating
\textbf{point values} through the chain $n\to I\to\delta\to E$ amplifies a
334\cmi{} input error by a factor of $1.3\times10^{6}$, because
$dI/d\delta\approx75\,000$\cmi. Propagating \textbf{intervals} through the same
chain does not amplify at all: each level's width is set by the spacing of its
own measured neighbours, not inherited from below. Monotone maps compose exactly
on intervals and catastrophically on points.

\section{Discussion}

\subsection{The thesis against the strongest premise}
We began expecting the lattice's order to constrain physics. It does not, and
the reason is not that the index was too coarse. Over the course of this work
the index was \textbf{improved to exactness} --- injective at eleven dimensions,
resolving core $J$, fine structure and the jK label --- and the predictive
boundary did not move by one cell.

The useful formulation is therefore: \textbf{an ordinal structure over
configuration space supplies indexing, and indexing is not prediction.} What the
lattice contributes to Part II is the well-definedness of ``same channel'' and
``interior on some axis'' --- real contributions, since without them the method
could not be stated --- and nothing further.

\subsection{Why energies and nothing else}
$T=\Zeff^{2}R/\nu^{2}$ is strictly monotone in every coordinate, has no zeros or
stationary points, varies by orders of magnitude across a series, and is
measured to $10^{-3}$\cmi. The bracket requires all four properties. We tested
the two nearest alternatives and each fails on a different one: dipole matrix
elements have zeros; fine-structure splittings are asymptotically constant. We
know of no other atomic observable with all four.

\subsection{Claims withdrawn}
\label{sec:withdrawn}
Stated explicitly so readers of earlier versions are not misled.
\begin{enumerate}[leftmargin=*,itemsep=1.5pt]
\item \textbf{$|\Le|=61\,453$ as an invariant.} A truncation artefact; all
derived region fractions likewise.
\item \textbf{``Zero invention'' as a substantive property.} The set measured at
0\% completion cost was defined by admissible chain constraints, and such a set
is closed by construction. The measurement was tautological. The genuinely
physical set has nonzero completion cost at every dimension.
\item \textbf{The physical set as a sublattice.} False; spin parity breaks join.
\item \textbf{Offset coordinates as a repair.} Faithful and injective but not a
lattice.
\item \textbf{Any closure-preserving $h$ depends on one coordinate.} True for
products of chains, false for \Lm{}: 148 of 335 counterexamples on a test
lattice. Lemmas 4.1--4.2 replace it and are what the argument needs.
\item \textbf{Ca\,\textsc{i} $3d^{2}\,^1$D$_2$ at $48\,800$--$49\,000$\cmi.}
Wrong by ${\approx}1000$\cmi; it places $^1$D above $^3$P, which for $d^{2}$
requires $C/B>5$~\cite{racah}. Revised unperturbed position
$47\,539$--$48\,213$\cmi.
\item \textbf{A point estimate for Ba\,\textsc{i} $5d^{2}\,^1$G$_4$.} Three
methods gave $24\,163$, $23\,852$ and $22\,094$\cmi, the last outside its own
bracket. Withheld-term validation fails by 1330\cmi. Only the bracket survives.
\item \textbf{Maximality without qualification.} No derived quantity raises the
dimension; independent quantities do.
\item \textbf{The $h$-axis impossibility as a general result.} It is an instance
of Theorem~4.3 --- $h$ was defined as a function of outputs.
\end{enumerate}

\subsection{Open questions}
Whether any atomic observable other than energy satisfies the four conditions of
\S15.2. Whether the multi-subshell coverage gap admits a representation
preserving closure --- Appendix~\ref{app:multi} shows the obvious route does
not. And whether the bracketing method transfers to non-atomic ordered systems,
which we have not attempted.

% =====================================================================
\appendix
\parttitle{Appendices}

\section{The multi-target extension and why it fails}
\label{app:multi}
The 86\% coverage ceiling is caused by configurations with two occupied
non-source subshells. The natural repair is a second target triple
$(e_2,f_2,g_2)$. It fails, instructively.

\textbf{Ordering the targets is free.} The canonicalisation $e_1\le e_2$ is
admissible --- $\varphi$ is the identity --- with zero failures.

\textbf{Conservation is not.} The requirement $g_1+g_2\le q$ is a bound on a
\textbf{sum}, hence outside \eqref{eq:admis}, and it breaks closure: $89\,864$
join failures in $979\,300$ sampled pairs, meets unaffected. The counterexample
is two lines --- $a=(g_1{=}1,g_2{=}0,q{=}1)$ and $b=(g_1{=}0,g_2{=}1,q{=}1)$
each conserve, and their join transfers two electrons having removed one.

\textbf{Dropping conservation restores closure} ($45\,690$ cells, zero failures)
and produces a lattice describing configurations that cannot exist.

\textbf{The conserving set is not a sublattice but is almost a lattice.} With
sufficient headroom, $133\,542$ of $133\,542$ interior pairs have a unique least
upper bound --- the join exists but is not componentwise max: $q$ rises to admit
both transfers.

\textbf{With two sources uniqueness fails too.} Conservation becomes
$g_1+g_2\le q_1+q_2$, a sum on both sides; 2.1\% of pairs acquire multiple
minimal upper bounds. The result is a graded bounded poset, neither a join- nor
a meet-semilattice.

\textbf{And the slot representation is not faithful.} Sorted pairs and multisets
are in bijection but the bijection is \textbf{not a lattice homomorphism}: for
$3d^{1}$ and $4p^{2}$, multiset union gives three transferred electrons and the
sorted-pair join gives two.

\section{Computational methods}
All computations in Python 3 with NumPy~\cite{numpy} and SciPy~\cite{scipy};
figures with Matplotlib~\cite{matplotlib}. Closure tests either exhaustive over
all pairs (stated where so) or on uniform random samples of 900--8\,000\,000
pairs with the sample size reported. Box searches enumerate all $2^{d}$ corners
of candidate base cells. Null models for the closure-as-existence tests are of
two kinds: uniform random subsets of the physical envelope matched in size, and
marginal-matched subsets drawn with probability proportional to the observed
frequency of each coordinate value, the latter constructed specifically to
absorb observational-completeness effects. Ritz fits use ordinary least squares
of $\delta$ against $1/n^{2}$, with $\sigma$ from the standard prediction error
including the leverage term
$s\sqrt{1+1/k+(x_0-\bar{x})^{2}/S_{xx}}$.

\section{Data provenance}
\label{app:data}
\begin{center}\small
\begin{tabular}{@{}p{62mm}p{78mm}@{}}
\toprule
source & used for\\
\midrule
NIST ASD ver.\ 5.12~\cite{nistasd} & Ca\,\textsc{i/ii}, Sr\,\textsc{ii},
  Kr\,\textsc{i/ii}, Sc\,\textsc{iii}, Ti\,\textsc{i} levels\\
Sansonetti \& Nave~\cite{sansonetti} & Sr\,\textsc{i} levels and transition
  probabilities\\
Curry~\cite{curry} & Ba\,\textsc{i/ii} levels, Land\'e factors\\
Martin \emph{et al.}~\cite{srd111} & ground configurations, ionisation energies\\
NIST Handbook~\cite{handbook} & H\,\textsc{i}, Mg\,\textsc{i} levels\\
Rafiq \emph{et al.}~\cite{rafiq} & Mg\,\textsc{i} $3snd\,^1$D$_2$ series\\
Cowley~\cite{cowley} & successive ionisation energies\\
\bottomrule
\end{tabular}
\end{center}
\vspace{2mm}
Levels are quoted as published; centres of gravity are $(2J+1)$-weighted.

\section{Errors found and corrected during this work}
\label{app:errors}
Recorded because the method by which they were caught is part of the result.
\begin{itemize}[leftmargin=*,itemsep=1.5pt]
\item \textbf{Ba\,\textsc{i} $6s11f$ mis-keyed as $6s8f$.} Detected by the
derived ordinal coordinate producing four order violations, all involving one
cell. The datum implies $\delta=-2.84$ as $n=8$, impossible, and $\delta=+0.16$
as $n=11$, correct for an $f$ orbital.
\item \textbf{Ca\,\textsc{i} $3d^{2}\,^1$D$_2$ position.} Detected by the
$d^{2}$ term-ordering constraint, which three fitted analyses had not applied.
\item \textbf{Regex parsing failure.} An early configuration parser read
\texttt{3d4s} as $3d^{4}$, silently dropping the second subshell.
\item \textbf{Bracket direction for decreasing sequences.} For a decreasing
property the tightest bracket takes the maximum over cells above in $n$ and the
minimum over cells below; the reverse returns the global extremes. The error
inflated a reported width from 11\% to 114\%.
\item \textbf{Pooled $\sigma$.} Quoting a single $\sigma$ across series with
genuinely different defect scatter (a $370\times$ range) produced 96.2\%
coverage at nominal $\pm1\sigma$ and was misread as heavy tails.
\item \textbf{Anchoring on a published estimate.} A point prediction for
Ba\,\textsc{i} $5d^{2}\,^1$G$_4$ was developed with a literature value visible
from the first computation; the reported range was 38\% narrower than the inputs
supported. This is why \S11.2 reports a bracket and no estimate, and why the
Kr\,\textsc{ii} prediction was committed in writing before the measurement was
requested.
\item \textbf{Isolated closure mistaken for joint closure.} The offset
coordinates of \S\ref{sec:offset} were each verified closed in their own plane
and asserted to be jointly closed. Exhaustive testing showed $222\,075$ join
failures.
\end{itemize}

% =====================================================================
\section*{Acknowledgement of computational assistance}
\addcontentsline{toc}{section}{Acknowledgement}
Substantial parts of the analysis --- the closure and dimension tests, the
ablation studies, the 295-cell verification, the failure diagnostics of Part
III, and several of the corrections in Appendix~\ref{app:errors} --- were
carried out in dialogue with an AI system (Claude, Anthropic), which executed
the computations and served as an adversarial check on intermediate conclusions.
The system does not hold authorship: authorship entails accountability it cannot
bear, since it will not persist to answer for the work. Its errors, several
recorded in Appendix~\ref{app:errors}, were caught by the same process that
caught the human ones. What can be said precisely is that no numerical claim in
this paper rests on unexamined output: each was recomputed, attacked, and in a
number of cases withdrawn.

% =====================================================================
\begin{thebibliography}{99}
\small

\bibitem{nistasd} A.~Kramida, Yu.~Ralchenko, J.~Reader and NIST ASD Team,
\emph{NIST Atomic Spectra Database}, version 5.12, National Institute of
Standards and Technology, Gaithersburg MD (2024).
\url{https://physics.nist.gov/asd}

\bibitem{sansonetti} J.~E. Sansonetti and G.~Nave, Wavelengths, transition
probabilities, and energy levels for the spectrum of neutral strontium (Sr\,I),
\emph{J. Phys. Chem. Ref. Data} \textbf{39}, 033103 (2010).

\bibitem{curry} J.~J. Curry, Compilation of wavelengths, energy levels, and
transition probabilities for Ba\,I and Ba\,II, \emph{J. Phys. Chem. Ref. Data}
\textbf{33}, 725 (2004).

\bibitem{srd111} W.~C. Martin, A.~Musgrove, S.~Kotochigova and J.~E.
Sansonetti, \emph{Ground Levels and Ionization Energies for the Neutral Atoms},
NIST Standard Reference Database 111.

\bibitem{handbook} J.~E. Sansonetti, W.~C. Martin and S.~L. Young,
\emph{Handbook of Basic Atomic Spectroscopic Data}, NIST Standard Reference
Database 108.

\bibitem{salomankr} E.~B. Saloman, Energy levels and observed spectral lines of
krypton, Kr\,I through Kr\,XXXVI, \emph{J. Phys. Chem. Ref. Data} \textbf{36},
215 (2007).

\bibitem{salomanti} E.~B. Saloman, Energy levels and observed spectral lines of
neutral and singly ionized titanium, Ti\,I and Ti\,II, \emph{J. Phys. Chem. Ref.
Data} \textbf{41}, 013101 (2012).

\bibitem{cowley} C.~R. Cowley, Successive ionization energies of the elements
$Z=1$--100, University of Michigan compilation.

\bibitem{rydberg} J.~R. Rydberg, Recherches sur la constitution des spectres
d'\'emission des \'el\'ements chimiques, \emph{Kongliga Svenska
Vetenskaps-Akademiens Handlingar} \textbf{23}, 1 (1889).

\bibitem{ritz} W.~Ritz, \"Uber die Spektren der Alkalien, \emph{Ann. Phys.}
\textbf{317}, 264 (1903).

\bibitem{seaton} M.~J. Seaton, Quantum defect theory, \emph{Rep. Prog. Phys.}
\textbf{46}, 167 (1983).

\bibitem{fano} U.~Fano, Quantum defect theory of $l$ uncoupling in H$_2$ as an
example of channel-interaction treatment, \emph{Phys. Rev. A} \textbf{2}, 353
(1970).

\bibitem{lufano} K.~T. Lu and U.~Fano, Graphic analysis of perturbed Rydberg
series, \emph{Phys. Rev. A} \textbf{2}, 81 (1970).

\bibitem{cookecromer} W.~E. Cooke and C.~L. Cromer, Multichannel quantum-defect
theory and an equivalent $N$-level system, \emph{Phys. Rev. A} \textbf{32}, 2725
(1985).

\bibitem{rafiq} M.~Rafiq, M.~A. Kalyar and M.~A. Baig, Measurement of
photoionization cross sections from the $3s\,nd\ ^1$D$_2$ Rydberg states of
magnesium, \emph{J. Phys. B} \textbf{40}, 3181 (2007).

\bibitem{racah} G.~Racah, Theory of complex spectra, I--IV, \emph{Phys. Rev.}
\textbf{61}, 186 (1942); \textbf{62}, 438 (1942); \textbf{63}, 367 (1943);
\textbf{76}, 1352 (1949).

\bibitem{cooper} J.~W. Cooper, Photoionization from outer atomic subshells: a
model study, \emph{Phys. Rev.} \textbf{128}, 681 (1962).

\bibitem{condonshortley} E.~U. Condon and G.~H. Shortley, \emph{The Theory of
Atomic Spectra}, Cambridge University Press (1935).

\bibitem{cowan} R.~D. Cowan, \emph{The Theory of Atomic Structure and Spectra},
University of California Press (1981).

\bibitem{birkhoff} G.~Birkhoff, \emph{Lattice Theory}, 3rd ed., American
Mathematical Society Colloquium Publications 25 (1967).

\bibitem{daveypriestley} B.~A. Davey and H.~A. Priestley, \emph{Introduction to
Lattices and Order}, 2nd ed., Cambridge University Press (2002).

\bibitem{dushnik} B.~Dushnik and E.~W. Miller, Partially ordered sets,
\emph{Amer. J. Math.} \textbf{63}, 600 (1941).

\bibitem{trotter} W.~T. Trotter, \emph{Combinatorics and Partially Ordered Sets:
Dimension Theory}, Johns Hopkins University Press (1992).

\bibitem{gratzer} G.~Gr\"atzer, \emph{Lattice Theory: Foundation},
Birkh\"auser (2011).

\bibitem{yannakakis} M.~Yannakakis, The complexity of the partial order
dimension problem, \emph{SIAM J. Alg. Disc. Meth.} \textbf{3}, 351 (1982).

\bibitem{ganterwille} B.~Ganter and R.~Wille, \emph{Formal Concept Analysis:
Mathematical Foundations}, Springer (1999).

\bibitem{madelung} E.~Madelung, \emph{Die mathematischen Hilfsmittel des
Physikers}, 3rd ed., Springer (1936).

\bibitem{klechkowski} V.~M. Klechkowski, \emph{The Distribution of Atomic
Electrons and the Rule of Successive Filling of $(n+l)$ Groups}, Atomizdat,
Moscow (1962).

\bibitem{scerri} E.~R. Scerri, \emph{The Periodic Table: Its Story and Its
Significance}, Oxford University Press (2007).

\bibitem{slater} J.~C. Slater, Atomic shielding constants, \emph{Phys. Rev.}
\textbf{36}, 57 (1930).

\bibitem{ostrovsky} V.~N. Ostrovsky, What and how physics contributes to
understanding the periodic law, \emph{Found. Chem.} \textbf{3}, 145 (2001).

\bibitem{numpy} C.~R. Harris \emph{et al.}, Array programming with NumPy,
\emph{Nature} \textbf{585}, 357 (2020).

\bibitem{scipy} P.~Virtanen \emph{et al.}, SciPy 1.0: fundamental algorithms for
scientific computing in Python, \emph{Nature Methods} \textbf{17}, 261 (2020).

\bibitem{matplotlib} J.~D. Hunter, Matplotlib: a 2D graphics environment,
\emph{Comput. Sci. Eng.} \textbf{9}, 90 (2007).

\end{thebibliography}

\vspace{4mm}
\noindent{\small\color{mut}\textbf{Note on citation practice.} Sources are cited
for data and for standard results used in the argument. Where a result is
standard and attributable to a textbook rather than a paper --- the dimension of
the Boolean lattice, the closure of a product of chains --- the textbook is
cited and no priority is claimed. No source is cited in support of a claim it
does not make; in particular the order-theoretic construction of \S2 is not, so
far as the author is aware, present in the literature, and the negative result
of Part III is offered as new.\par}

\end{document}
